\section{Lazy integer encoding}\label{part:encoding}

The conflict analysis of \cref{sec:learning} resolved over Boolean bound
literals such as $\bl{x_5 \ge 6}$ and learned clauses written in them, taking for
granted that those Booleans exist. This section describes where they come from:
the part of CP-SAT that reconciles an integer variable with a huge domain with a
CDCL engine that ultimately reasons in Boolean literals.

The lazy-reason machinery of the preceding sections only works because the
Boolean layer it feeds is itself built lazily: an integer variable's encoding
into Booleans is created on demand, the same way its reasons are. That is why it
follows learning directly rather than waiting among the later accelerators.

The encoding below puts one Boolean on each bound $\bl{x \ge v}$, which looks like a \emph{unary} encoding: a variable over a million
values would seem to need a million Booleans, the very blow-up a binary encoding
exists to avoid. The resolution is the word \emph{lazy}. CP-SAT never lays down
that unary chain in full; it mints a bound literal only at the moment some
mechanism actually refers to it, so a wide integer typically contributes a
handful of Booleans rather than its whole range. The encoding is unary in
\emph{form} but sparse in \emph{practice}, and the rest of this section is about
how that sparsity is achieved and what it costs elsewhere.

It also helps to be precise about what kind of mechanism this is. Lazy clause
generation (\cref{sec:learning}) is a \textbf{correctness} mechanism:
reasons are what make learned clauses sound deductions. Lazy \emph{encoding}, by
contrast, is an \textbf{efficiency} mechanism: CP-SAT would stay correct if it
encoded every bound up front. Whether that would even hurt depends on the
domains involved. Many models are dominated by Boolean variables, and a variable
with a small domain costs only a few bound literals; there, eager encoding
would do no real harm. The cost concentrates on the variables whose domains are
genuinely wide. The objective is the usual one (a weighted sum can range over
millions), and CP-SAT mints further wide-domain integers of its own when it
breaks a long linear constraint into a tree of partial sums. Encode every
bound of those up front and the model turns slow and memory-hungry.

\subsection{Minting Booleans on demand}\label{sec:encoding}
% Sources (VERIFIED): ortools/sat/integer.h, ortools/sat/integer.cc
%   IntegerEncoder; GetOrCreateAssociatedLiteral() (order/"list" encoding,
%   adds implications to neighbor literals; comment cites Feydy & Stuckey,
%   "Lazy Clause Generation Reengineered", CP 2009);
%   GetOrCreateLiteralAssociatedToEquality(); FullyEncodeVariable() (eager path);
%   Canonicalize() (folds domains with holes); NegationOf() shares the encoding;
%   num_created_variables_ counts created Booleans.
%   ortools/sat/integer_search.cc -- GetDecisionLiteral() calls
%     GetOrCreateAssociatedLiteral when branching (one of the two creation sites).
%   ortools/sat/integer_resolution.cc -- the other creation site (1-UIP).
%   ortools/sat/cp_model_loader.cc -- ExtractEncoding() creates value literals
%     (x == v) up front during loading/presolve (value encoding is STATIC).

\textbf{Why eager encoding is hopeless.} The naive bridge would give every
bound its own Boolean up front. A variable $x \in [0, 10^6]$ would spawn a
million literals
\[
  \bl{x \ge 1},\quad
  \bl{x \ge 2},\quad
  \ldots
\]
before search even starts. This is a non-starter in both memory and propagation
cost, and the overwhelming majority of those literals would never be touched.
CP-SAT therefore does not pre-encode all bounds. The integer trail reasons
about bounds \emph{directly}, as native integer facts
(\cref{sec:bounds}), and Booleans are created only where the integer and
Boolean worlds actually need to meet.

\textbf{The order encoding.} When that bridge is built, it uses the
\emph{order encoding}, also called the ``list encoding.'' CP-SAT follows Feydy
and Stuckey's \emph{``Lazy Clause Generation Reengineered''} (CP
2009)~\cite{feydy2009lcg}.
% The code comment at integer.h:106 attributes the order encoding to this paper.
For a variable $x$, the bound
literals have the form $\bl{x \ge v}$, and they are linked by
the obvious monotone implications:
\[
  \bl{x \ge 6} \;\Rightarrow\; \bl{x \ge 5}
  \;\Rightarrow\; \bl{x \ge 4} \;\Rightarrow\; \cdots .
\]
A single such Boolean also represents the corresponding upper-bound fact,
because
\[
  \bl{x \le v}
  \equiv
  \lnot\bl{x \ge v{+}1} .
\]
In the implementation, the upper-bound side is handled by sharing the encoding
with \code{NegationOf}$(x)$, so $x$ and its negation do not duplicate Booleans.

\textbf{On-demand creation.} The workhorse is
\srcla{integer.h}{193}{GetOrCreateAssociatedLiteral}$(x \ge v)$. It
returns the Boolean $\bl{x \ge v}$. If that literal does
not yet exist, it mints a fresh Boolean variable and wires in the implications to
the already encoded neighboring bounds, the next bound below and the
next bound above. This preserves the monotone chain.

This is the moment alluded to in the conflict trace of \cref{sec:trace}.
When 1-UIP resolution needs $\bl{x \ge v}$ as a clause literal, this call
materializes it. The implication links keep it consistent
with its neighbors, and from then on the SAT engine treats it like any other
Boolean literal. In the basic search loop, two moments dominate bound-literal
creation: \emph{branching} on an integer bound mints the Boolean for that
bound (\srcla{integer\_search.cc}{1426}{GetDecisionLiteral}), and
\emph{integer conflict resolution} mints the Boolean it needs for the 1-UIP
(\cref{sec:trace}). Other mechanisms call the same
\code{GetOrCreateAssociatedLiteral} on demand as well, including core-guided
optimization when forming assumption literals, the feasibility pump, and some
global propagators. Ordinary bound propagation, however, stays in integer-bound
space and creates no Booleans. That is what keeps the Boolean skeleton small.

\textbf{Equality literals.} Some constraints, for example a table constraint,
\srcla{all\_different.h}{52}{AllDifferent}, or anything that reasons about exact
values, want $\bl{x = v}$, not just bounds.
\srcla{integer.h}{194}{GetOrCreateLiteralAssociatedToEquality}$(x, v)$
provides it, based on the identity
\[
  \bl{x = v} \;\Longleftrightarrow\;
  \bl{x \ge v} \;\wedge\; \bl{x \le v} .
\]
A constraint may even ask for the whole domain at once via
\srcla{integer.h}{144}{FullyEncodeVariable}. This is the deliberate
\emph{eager} path, used precisely when a propagator is going to need every value
anyway, so laziness would only add overhead. Full encoding is the exception
requested by a constraint; lazy bound encoding is the default for everything
else.

There is an important asymmetry between the two encodings. The bound literals
$\bl{x \ge v}$ of the order encoding are predominantly \emph{dynamic}: they
are minted during branching and conflict resolution as just described. Value
literals $\bl{x = v}$ are predominantly \emph{static}: the
equalities a model needs are detected and created up front during
presolve/loading (\srcla{cp\_model\_loader.cc}{352}{ExtractEncoding}), rather
than on demand. That said,
\code{GetOrCreateLiteralAssociatedToEquality} can still mint one lazily if a
constraint later asks for a value literal that loading did not foresee. The
developers' rule of thumb behind this split is blunt: if a model genuinely needs
per-value Booleans for a variable, that variable probably should not have been a
wide integer in the first place.

\textbf{Domains with holes.} Real domains are not always intervals. If
$x \in [1,2] \cup [5,6]$, then ``$x \le 3$'' and ``$x \le 4$'' mean the same
thing as ``$x \le 2$,'' and ``$x \ge 3$'' means ``$x \ge 5$.'' The encoder's
\srcla{integer.h}{177}{Canonicalize} step folds each requested bound onto
the canonical bound for that domain before looking it up. Equivalent bounds
therefore share one Boolean instead of spawning redundant ones.

\textbf{The payoff.} The Boolean skeleton that the CDCL engine actually sees
stays small: it contains only the bounds and equalities that branching,
conflict learning, or a constraint-specific encoding genuinely reached. A
variable over a million values typically contributes a handful of literals, not
a million, and the encoder tracks exactly how many Booleans it had to create
(\code{num\_created\_variables\_}). Both layers are lazy: the reasons
(\cref{sec:learning}) and the Boolean variables they are written in.
Strong integer propagation does
the bulk of the reasoning natively over bounds, and the Boolean encoding
materializes only at the thin interface where the learning machinery needs to
grip.

% Source (VERIFIED): "Search stats" column "Bools" = FormatCounter(r.num_booleans())
%   in ortools/sat/stat_tables.cc:90. The response field "booleans:"
%   (cp_model_solver.cc:711) is one worker's SatSolver::NumVariables(), not a
%   portfolio figure -- different workers mint different amounts (see the portfolio
%   box in the parallelization section).
\begin{platypuslog}[htb]
The size of the lazily minted Boolean skeleton is reported per worker as the
\code{Bools} column of the \textsf{Search stats} table; each worker mints its own,
so the counts differ across rows. For a model with wide integer domains it usually
stays orders of magnitude below the number of representable values: the measurable
payoff of minting bound literals on demand rather than up front.
\end{platypuslog}

\subsection{A fixed skeleton need not be a solution}\label{sec:completion}
% Sources (VERIFIED):
%   ortools/sat/integer_search.cc:422 SatSolverHeuristic -- all_assigned =
%     trail->Index() == sat_solver->NumVariables(); returns no decision when true.
%   The completion heuristic that closes this gap (AtMinValue, bisection, LP
%   value; completion-last policy) is assembled in sec:search.

The sparse skeleton has a consequence the foundational search of
\cref{sec:complete} sidestepped. Because the skeleton holds only a handful
of bounds, all of them can be assigned while an integer still spans a range:
$\bl{x \ge 3}$ true and $\bl{x \ge 8}$ false
leave $x \in [3,7]$ with every \emph{existing} literal decided. The gap is then a
\emph{missing} literal, not an unassigned one. Even a complete Boolean assignment
can leave $x$ genuinely ranged, because no bound between $3$ and $7$ was ever
minted. Under a full encoding this could not arise, since deciding every
bound pins the value; the sparse skeleton is exactly what pulls the two
tests apart.

Closing this gap is the job of \emph{integer completion}, the second half of the
search machinery in \cref{sec:search}: once the Boolean skeleton is
settled but some integer is still ranged, the solver mints a fresh bound for
that integer and branches on it, repeating until every variable is pinned to a
single value.
